
\documentclass[11pt,a4paper]{article}

% --- Packages ---
\usepackage[utf8]{inputenc}           % Fixed encoding package name
\usepackage{amsmath, amssymb, amsthm} % Mathematical typesetting
\usepackage{booktabs}                 % Professional-grade table lines
\usepackage{geometry}                 % Page layout definitions
\usepackage{microtype}                % Improves justification and spacing
\usepackage{titlesec}                 % Custom section styling
\usepackage{xcolor}                   % Box backgrounds and textcoloring

% --- Layout Configuration ---
\geometry{
    top=1in,
    bottom=1in,
    left=1in,
    right=1in
}

% --- Custom Styling Elements ---
\titleformat{\section}{\large\bfseries}{\thesection}{1em}{}[{\titlerule[0.5pt]}]
\titleformat{\subsection}{\normalsize\bfseries}{\thesubsection}{1em}{}

% Define a styled blockquote environment for deep highlights
\newenvironment{technicalnote}{%
    \medskip
    \noindent
    \color{black!90}%
    \vrule width 3pt \hspace{10pt}%
    \begin{minipage}{\dimexpr\textwidth-15pt}%
    \itshape
}{%
    \end{minipage}%
    \medskip
}

% --- Document Metadata ---
\title{\textbf{Empirical Measurement of Temporal Jacobian Instability in LSTM-PPO Agents Under Partial Observability }}

\begin{document}

\maketitle

\section{Introduction: }
Reinforcement Learning (RL) agents frequently fail during training. While these failures are commonly attributed to "bad hyperparameters", RL agents differ significantly from supervised learning agents; they are notoriously difficult to train, and diagnosing failures is inherently challenging. The underlying cause is often numerical: unstable gradient dynamics, ill-conditioned weight matrices, and representation drift that compound over time. 

\section{Background: }
\subsection{The Hard Problem of Stability Verification}
Proving or certifying global analytical stability in Rl is mathematically impossible due to three compounding factors: 

\begin{itemize}
    \item \textbf{Highly Non-linear Policy and Environment :} Deep neural networks are a massive composition of non-linear functions (e.g., the hidden state of an LSTM transitions through time as a non-linear dynamical system). Because of this non-linearity, the system's sensitivity to inputs is non-uniform. A small change in an environment's state vector might cause a negligible change in the network's output action in one region, but a catastrophic, exploding change in another, which is a manifestation of the local curvature phenomenon. 
    \item \textbf{Closed Feedback Loops: } Real-world environments are typically governed by non-linear differential equations. When a non-linear policy is integrated into a non-linear environment, they form a closed feedback loop. THe policy outputs an acrtion based on its current weights, and the environment accepts this action to compute the next state via a non-linear transition function. Because both systems are non-linear, errors do not merely add up; they multiply exponentially. This triggeres extreme sensitivity to initial conditions. 
    \item \textbf{Dynamically Changing Policy: } IN RL, the system transforms while it runs. As the agent interacts with the environment, gradient descent updates the weights, meaning the policy is a moving target. the optimization landscape for RL agents is a deforming dynamical system. The  basins of attraction and loss cliffs are continously shifting, expanding and swallowing parameters. Because the policy is constantly moving, the agent encounters entirely novel states, leading to unpredictable errors at the input level. These errors propagate through the network, causing massive gradient spikes that completely alter the weights and can destroy hours of stable learning. 
\end{itemize}

Without explicit regulation of the gradients, RL agents exhibit dangerous behaviors in the long run. Performance frequently starts to deteriorate due to large gradient variance and high sensitivity to step size, eventually losing previous gains and behaving erratically. This occurs because the Lipschitz gains grow unboundedly without regulation. 
The policy network's partial gradients, $\partial j \pi i(x)$ (which measures how sensitive each output is to each input) are unconstrained during standard training. AS training progresses, gradient updates push these partial derivatieves larger and larger this manifests as the Lipschitz constant of the policy growing unboundedly over iterations when no gradient regulation is applied. 
Conceptually, the policy network becomes progressively more "jumpy" reacting more dramatically to smaller shanges in its observation. If this sensitivity happens to increase rewards along a specific trajectory, nothing in the standard training process penalizes it. While the agent remains in forgiving situations, this instability is masked. However, once training pushes the agent into a slightly different, trickier situations, a tiny state perturbation causes a massive overreaction. This overreaction alters the subsequent states, compunding into a feedback spiral that causes performance to collapse completely. THis is the same bifurcation phenomenon observed in recurrent nueral network (RNN) training dynamics. 

\subsection{Why LSTM-PPO}
Standard feedforward RL algorithms rely entirely on the Markov Property, which assumes the environment's future depends only on its present state, not its past history. When an agent operates in a partially observable environment, the Markov property is broken from the agent's perspective. To make accurate predictions and decisions, the agent requires the history of past observations, necessitating recurrence (e.g., an LSTM). 
The recurrenct network maintains an internal, evolving hidden state $h_t$ that compresses the historical trajectory. The policy then maps both the current observation and this memory state to an action, allowing the network to effectively model complex, higher-order dynamical systems. 



\section{Mathematical Pathalogies Introduced by Reccurrence}
When an LSTM learns via Backpropagation Through Time (BPTT), it propagates an error signal backward through every time step. At each timestep, this signal is multiplied by a weight matrix. Repeatedly multiplying a vector by a matrix causes the signal to either grow exponentially or decay to zero, depending heavilt on the spectral radius of the matrix. 

\subsection{The fragile Zone $\&$ Spectral Radius}
The spectral radius $\rho(W)$ of a matrix $W$ is the magnitude of its largest eigenvalue: 
\begin{equation}
    \rho(W) = \max_i |\lambda_i|
\end{equation}
It dictates the worst-case growth or shrinkage rate under repeated matrix multiplication. 
\begin{itemize}
    \item  $\rho(W) > 1$: Repeated multiplication explodes. 
    \item $\rho(W) < 1$: Repeated multiplication vanishes.
    \item $\rho(W) = 1$: The fragile zone, where information persists but stability is marginal. 

\end{itemize}
If a system needs to retain long-term memory, its internal memory matrix requires eigencalues very close to 1. Consequently, the system becomes extremely sensitive to tiny changes; long memory and absolute numerical stability are fundamentally in tension within standard recurrenct architectures. 
This structural tension explains why optimization-level constraints alone cannot prevent internal numerical collapse. For instance, algorithms like Trust Region Policy Optimization (TRPO) use a localized safety boundary (via a KL-divergence constraint) to guarantee that the agent's visible behavior and performance imprve smoothly. However, this safety restriction only applies to external actions (policy space). Underneath these stable outputs, the internal gradient updates, recurrent Jacobians, and hidden memory states are still free to fluctuate wildly. This disconnect allows deep numerical pathalogies to grow completely unchecked, which is exactly why measurinf spectral radius variance is necessary to capture the true internal health of the network. 

\subsection{Spectral Gap, Non-Normality, and Bifurcations}
The spectral gap measures the distance between the spectral radius and the stability boundary $(1-\rho(W))$, or alternatively, the distance between the largest and second-largest eigenvalues. It quantifies proximity to the edge of stability: a small spectral gap indicates the system is in a dangerous, highly sensitive regime. 

The Spectral SSM framework complements this backward pass perspective by analyzing the forward pass of a linear dynamical system, $x_t = Ax_{t-1} +Bu_t$. Unrolling this system reveals that the state at step $t$ depends on repeated powers of the dynamical matrix,$A^t$. If all eigenvalues of $A$ are bounded away from 1 by a spectral gap $\delta$ (such that $\rho(A) <= 1 - \delta$), the historical information decays at a rate of $(1-\delta)^t \approx e^{-\delta t} $. This establishes that the effective memory horizon of the network is on the order of $1/\delta$. Consequently, as $\delta -> 0$ to allow for long-term memory retention, the system inherently appoches marginal stability, mirroring the exact numerical fragility encountered during backpropagation.

In recurrent neural networks, weight matrices are typically non-normal ($W^tW \not = WW^t$). Non-normal matrices are notorious because their eigenvalues can be highly sensitive to small perturbations. When a matrix has a small spectral gap and is highly non-normal, its eigenvectors can become nearly parallel.

When this occurs, the condition number of the eigenvector matrix becomes massive. A high condition number implies that an infinitesimally small update to the weights can cause a disproportionately large, discontinuous leap in the eigenvalues. In dynamical systems, a bifurcation often occurs exactly when the spectral gap shrinks to zero or an eigenvalue crosses a critical threshold. As the network approaches a loss cliff, the spectral gap closes, the system becomes heavily ill-conditioned, and a tiny gradient step can trigger a global structural change—a bifurcation causing eigenvalues to scatter and the spectral radius to shoot up instantaneously.

\subsection{The Local Curvature Phenomenon}
The local curvature phenomenon refers to the rapid and extreme increase in the curvature of the error surface in RNNs, specifically in directions along which gradients explode. In general nonlinear RNNs, there exists a direction such that if the gradient explodes along it, the curvature along that same direction also increases exponentially. 

This explosion of curvature creates geometric structures in the error surface known as highcurvature wells or cliff-like structures. These walls are perpendicular to the direction of the exploding gradient. Around these walls, the error surface is characterized by a steep cliff followed by a smoother, low-curvature region.

When a standard gradient descent step encounters one of these high-curvature walls, the calculated gradient magnitude becomes massive. This forces the optimization step to jump violently across the valley, moving perpendicular to the wall. This jump can throw the parameters very far into an unrelated region of the parameter space, disrupting the learning process entirely. The higher variance observed in recurrent RL is a direct geometric manifestation of these underlying dynamical events, such as basin crossing

\subsection{Why Spectral Radius variance is Higher under Recurrence}

To learn long-term dependecies, the spectral radius of the recurrent weight matrix must often hover near or above 1. As the spectral radius increases to accommodate longer sequences, the modelenters complex dynamical regimes containing multiple attractors and intricate boundaries between basins of attraction. 

In the reccurrent case with long sequences, variance is amplified due to: 

\begin{itemize}
    \item \textbf{Amplified Sensitivity:} The factors transporting error through time are products of Jacobian matrices. When the spectral radius exceeds unity, any small modification to the dynamical system is amplified exponentially over time. 
    \item \textbf{Constant reorganization:} Because the model must balance between vanishing and exploding gradients, the effective stability of the system and consequently the observed variance of its spectral radius and gradients fluctuates wildly as it interacts with high curvature walls. 
    \item \textbf{Bifurcation Effects:} As training progresses and recurrent weights evolve, the model frequently transitions between stable and unstable regimes via local crossings of global bifurcation boundaries, resulting in severe variance spikes. 
\end{itemize}
\section{Identified Research Gaps}
Despite the rigorous foundations established in literature like Pascanu's work, several major gaps persist in current RL research:
\begin{itemize}
    \item \textbf{Lack of Empirical RL Validation:} Pascanu's theoretical framework has never been empirically measured in an active Reinforcement Learning setting.
    \item \textbf{Observability Limitations:} Stability analysis in RL has historically only been performed in fully observable settings.
    \item \textbf{Missing Baseline Comparisons:} No direct numerical comparison between LSTM and feedforward agents on gradient stability metrics currently exists.
    \item \textbf{The SSM Translation Gap:} The spectral gap problem identified in spectral State-Space Model (SSM) papers has no empirical counterpart or validation in RL.
    \item \textbf{Isolated Analysis Spaces:} Policy-space stability and gradient-space stability have never been studied simultaneously in the same experiment.
    \item \textbf{Uncharacterized POMDP Interactions:} The interaction between partial observability and recurrent gradient instability remains completely uncharacterized.
    \item \textbf{No Standardization:} No reproducible measurement framework exists for parsing this specific problem.
\end{itemize}

\section{Objectives of the study: }
To bridge these gaps, this study will target four primary objectives:
\begin{enumerate}
    \item \textbf{Establish a standardized measurement framework} for evaluating numerical stability in recurrent RL.
    \item \textbf{Empirically test the direct effect of recurrence} on gradient stability under strictly controlled partial observability.
    \item \textbf{Connect empirical numerical instabilities} back to established, classical theoretical predictions.
    \item \textbf{Provide a diagnostic baseline} for recurrent RL numerical behavior to inform future algorithmic design and stabilization techniques.
\end{enumerate}

\section{Central Hypothesis: Why This Study Measures Spectral Radius Variance}
This study tracks the spectral radius over training and quantifies its variance. The spectral gap provides the structural explanation for why these values fluctuate. In this context, \textbf{spectral radius variance serves as the measurable empirical fingerprint of the underlying numerical pathology}: the bifurcations, the cliff crossings, the ill-conditioning, and the fundamental tension between long-term memory and numerical stability.

A feedforward network utilizing a finite history window structurally avoids these phenomena: its gradient product is bounded by a fixed number of steps $k$, it contains no recurrent eigenvalues drifting toward 1, and its spectral gap never closes as a function of sequence length.

\textbf{Central Hypothesis:} Spectral radius variance will be measurably and significantly higher for the LSTM than for the feedforward baseline---not due to the choice of optimizer or environment, but because of the mathematical conditions that recurrence inherently demands.


\section{Methodology}

\subsection{Environment}

We use a single, minimal partially observable GridWorld. The agent observes only a small window around its current position, defined by a fixed observation radius $r$; everything outside that radius is hidden. The degree of partial observability is held constant across all experiments. The environment is kept deliberately simple because it is not the variable of interest: it is a controlled vehicle for generating trajectories with meaningful partial observability, where the missing information is genuinely decision-relevant and cannot be recovered from a single observation.

The reward structure is sparse: a positive reward for reaching the goal and a small negative reward per step to discourage random wandering. Sparse reward is important because it forces both agents to use their memory mechanism to succeed, ensuring that the partial observability does real work rather than being trivially bypassed.

\paragraph{Grid connectivity guarantee.}
After placing the agent and goal on randomly selected empty interior cells, a breadth-first search (BFS) is run from the agent's position to verify that the goal is reachable. If the goal is unreachable under the current wall configuration, the grid is resampled until a valid layout is found. This ensures that no episode is silently unsolvable, regardless of the random seed.

\paragraph{Environment validation.}
Before any experiments, the environment is validated using its \texttt{full\_obs=True} mode, which returns the complete flattened grid as the observation instead of the local window. A fully observable agent trained under identical conditions should solve the task significantly faster than a partially observable one. If no performance gap exists, the observability masking is insufficient and the environment parameters must be adjusted.

\subsection{Agents}

Both agents face the same partial observability problem but solve it through architecturally different memory mechanisms. Everything else (environment, optimizer, hyperparameters, training duration, and random seeds) is held identical between them. Both agents operate under partial observability throughout all experiments; there is no fully observable training stage.

\paragraph{Agent 1: LSTM.}
A single-layer LSTM-based recurrent agent. At each step it receives the current observation $o_t$ and updates its hidden state:
\begin{equation}
  h_t,\, c_t = \mathrm{LSTM}\!\left(o_t,\, h_{t-1},\, c_{t-1}\right).
\end{equation}
The policy and value function are conditioned on $h_t$. The hidden state is reset to zero at the start of each episode and carries information across the full episode length. During training, gradients flow backward through every time step via backpropagation through time (BPTT), producing the temporal Jacobian product that the hypothesis concerns. The agent is trained with Proximal Policy Optimization (PPO).

\paragraph{Agent 2: Feedforward with finite history window.}
A feedforward network that receives the last $k$ observations concatenated into a single input vector,
\begin{equation}
  \tilde{o}_t = \left[\, o_{t-k+1},\, \ldots,\, o_t \,\right].
\end{equation}
There is no hidden state and no recurrence. The gradient product is bounded by $k$ steps by construction: the network's Jacobian never accumulates beyond that window. At episode start, the buffer is padded with zeros. The agent is trained with PPO under identical hyperparameters.

\paragraph{Choice of $k$.}
The history window is set to $k = \lfloor n/2 \rfloor$, where $n$ is the side length of the grid, giving $k = 5$ for the $10 \times 10$ grid and $k = 7$ for the $15 \times 15$ grid. This is large enough to capture recent trajectory information and give the agent a fair opportunity at the task, but not so large that it approximates unlimited memory. The choice is grounded in the grid geometry: half the side length is a reasonable lower bound on the number of steps needed to navigate from a random start toward the goal.

\subsection{Primary Measurement: Spectral Radius of the Temporal Jacobian Product}

At regular intervals during training (every $N$ gradient updates), the spectral radius of the temporal Jacobian product is computed for both agents on a held-out rollout:
\begin{equation}
  \rho_t \;=\; \rho\!\left(\prod_{i=k}^{T} \frac{\partial h_i}{\partial h_{i-1}}\right),
  \qquad
  \rho(A) = \max_j \left|\lambda_j(A)\right|,
\end{equation}
where $\lambda_j(A)$ denotes the eigenvalues of $A$.

For the LSTM agent, this product spans the full episode. For the feedforward agent, because there is no recurrent hidden state, the gradient product is bounded by $k$ steps by construction and its spectral radius is expected to be near zero. It therefore serves as a structural lower bound rather than a symmetrical measurement. The comparison is valid: the feedforward baseline establishes what gradient propagation looks like when recurrence is removed, but the two quantities are not computed identically and should not be interpreted as measuring the same object.

\paragraph{Implementation.}
Because \texttt{nn.LSTM} fuses the full sequence into a single kernel and does not expose per-step hidden states to backward hooks, the LSTM agent is implemented as a manual step-by-step loop in Python. At each timestep $t$, the forward pass produces $h_t$ explicitly, and \texttt{torch.autograd.grad} is used to compute $\partial h_t / \partial h_{t-1}$ with \texttt{retain\_graph=True}. The per-step Jacobians are accumulated into the product numerically. This computation is performed in a separate forward pass on a held-out rollout and is completely isolated from the optimizer, so it does not interfere with training gradients.

The result is two time series of spectral radius values over the full course of training, one per agent.

\subsection{Secondary Measurements}

The following quantities are logged alongside the primary measurement at the same intervals.

\paragraph{Gradient norm per update.}
The overall magnitude of the gradient used for each weight update. Spikes here should correlate with spectral radius spikes if the theoretical predictions hold.

\paragraph{Reward per episode.}
The agent's actual learning performance. If spectral radius variance spikes coincide with reward drops, this directly links the numerical instability being measured to behavioral collapse, providing the empirical connection to the findings of Jin and Lavaei.

\paragraph{Hidden state drift (LSTM only).}
The $\ell_2$ distance between consecutive hidden states, averaged over an episode:
\begin{equation}
  \frac{1}{T} \sum_{t=1}^{T} \left\lVert h_t - h_{t-1} \right\rVert_2.
\end{equation}
This measures how much the agent's internal memory fluctuates, and is a behavioral correlate of the eigenvalue instability that the spectral radius captures mathematically.

\paragraph{Condition number of the recurrent weight matrix $W_h$ (LSTM only).}
Tracked at checkpoints throughout training:
\begin{equation}
  \kappa(W_h) = \frac{\sigma_{\max}(W_h)}{\sigma_{\min}(W_h)}.
\end{equation}
A rising condition number should precede spectral radius spikes, since a high condition number is what makes the eigenvalues sensitive to tiny perturbations.

\subsection{What Confirms or Rejects the Hypothesis}

The primary comparison is the variance of the two spectral radius time series across the full training run. Beyond variance, the full distribution is examined, specifically whether the LSTM time series exhibits heavier tails or sudden large spikes, which would be the empirical signature of the bifurcation events and high-curvature wall crossings described in the conceptual background.

A secondary analysis checks whether spectral radius spikes correlate temporally with reward drops and gradient norm spikes. If they do, this establishes that the internal numerical instability is not merely a mathematical artifact but a real driver of behavioral collapse, connecting the measurement directly to the empirical findings of Jin and Lavaei in a recurrent, partially observable setting they did not study.

All experiments are run across a minimum of five random seeds per agent. The comparison is between distributions, not individual runs, ensuring statistical reliability.

\paragraph{Hypothesis supported if:}
the LSTM spectral radius variance is statistically significantly higher than the feedforward baseline across seeds, with the LSTM distribution showing heavier tails and occasional large spikes that are absent from the feedforward distribution.

\paragraph{Hypothesis rejected if:}
the two distributions are statistically similar. This would itself be a meaningful finding, suggesting that LSTM gradient dynamics in reinforcement learning are more robust in practice than the theoretical framework predicts.

\subsection{Why This Design Tests Exactly the Right Thing}

Both agents use PPO, so policy-space stability is controlled equally. Any difference in spectral radius variance cannot be attributed to the optimizer; it can only come from the architecture. This directly addresses Gap~5 from the research gaps: it is the first experiment that isolates architectural gradient-space instability from optimization-level instability, by holding the optimizer constant and varying only the memory mechanism.

The partially observable environment directly addresses Gap~6, the interaction between partial observability and recurrent gradient instability. The feedforward agent's bounded gradient product directly addresses Gap~3, the first direct numerical comparison between LSTM and feedforward agents on gradient stability metrics. The spectral radius measurement itself directly addresses Gap~1, the first empirical measurement of the temporal Jacobian spectral radius in an active reinforcement learning training loop.
\end{document}
